% GENERATED FILE. Edit canonical lessons, then run npm run generate:books.
% canonical-source-hash: fnv1a64:10b39303c3a92a16
% canonical-lessons: JA-C153-gaishutsu, JA-C153-kitaku, JA-C153-wasuremono, JA-C153-nakusu, JA-C153-maniau

\chapter{Out, Back Home and Things Left Behind}
\label{ch:ja-out-back-home-and-things-left-behind}
\hlchaptermodality{153}

\begin{chapteropening}
\textbf{By the end of this chapter:} \emph{I can talk about going out, coming home and losing things.}

Complete the last lesson of chapter 153: I can talk about going out, coming home and losing things.
\end{chapteropening}

\section[\texorpdfstring{\ja{がいしゅつ} (gaishutsu)}{がいしゅつ (gaishutsu)}]{\ja{がいしゅつ} (gaishutsu) — going out; being out of the office}

\label{lesson:JA-C153-gaishutsu}



\begin{quote}
Before the new one: say the Japanese for to forget, then the Japanese for to be late.
\end{quote}

\subsection*{You'll want to know: \ja{がいしゅつ}}
\textbf{\ja{がいしゅつ}} — \emph{gaishutsu} — \textquotedblleft{}going out; being out of the office\textquotedblright{}.

\begin{quote}
\textbf{\ja{ちちは} \ja{いま} \ja{がいしゅつちゅうです。}} — \emph{chichi wa ima gaishutsuchū desu} — \textquotedblleft{}My father is out right now.\textquotedblright{}
\end{quote}

\begin{cousinweb}
\textbf{\ja{がいしゅつ}} (\emph{gaishutsu}) is Sino-Japanese (kango): two Chinese characters read with their on'yomi, readings borrowed from Chinese. They are \emph{gai}, \textquotedblleft{}outside\textquotedblright{} (Middle Chinese \emph{ngwajH}) and \emph{shutsu}, \textquotedblleft{}go out\textquotedblright{} (\emph{tsyhwit}) (KANJIDIC2; Middle Chinese in Baxter's transcription).
\end{cousinweb}

\begin{grammarlens}[title={What you've built}]
One more word for everyday life at A2.
\end{grammarlens}

\subsection*{Your turn}
\begin{itemize}
\raggedright
  \item \emph{Say it:} \emph{gaishutsu}
  \item \emph{Say it:} \emph{gaishutsu}, once more
  \item \emph{Say it:} \emph{chichi wa ima gaishutsuchū desu} — My father is out right now.
  \item \emph{Recall:} say the Japanese for to forget, then the Japanese for to be late, then say \emph{gaishutsu} again
\end{itemize}

\begin{tcolorbox}[breakable,colback=teal!4,colframe=teal!35!black,title={Before you move on}]
What does \ja{がいしゅつ} mean? (\textquotedblleft{}Going out; being out of the office\textquotedblright{}.) Say it once more.
\end{tcolorbox}
\section[\texorpdfstring{\ja{きたく} (kitaku)}{きたく (kitaku)}]{\ja{きたく} (kitaku) — coming home; getting home}

\label{lesson:JA-C153-kitaku}



\begin{quote}
Before the new one: say the Japanese for to be late, then the Japanese for going out.
\end{quote}

\subsection*{You'll want to know: \ja{きたく}}
\textbf{\ja{きたく}} — \emph{kitaku} — \textquotedblleft{}coming home; getting home\textquotedblright{}.

\begin{quote}
\textbf{\ja{ちちの} \ja{きたくは} \ja{ゆうがたです。}} — \emph{chichi no kitaku wa yūgata desu} — \textquotedblleft{}My father gets home in the evening.\textquotedblright{}
\end{quote}

\begin{cousinweb}
\textbf{\ja{きたく}} (\emph{kitaku}) is Sino-Japanese (kango): two Chinese characters read with their on'yomi, readings borrowed from Chinese. They are \emph{ki}, \textquotedblleft{}return home\textquotedblright{} (Middle Chinese \emph{kjw+j}) and \emph{taku}, \textquotedblleft{}house, home\textquotedblright{} (\emph{draek}) (KANJIDIC2; Middle Chinese in Baxter's transcription).
\end{cousinweb}

\begin{grammarlens}[title={What you've built}]
One more word for everyday life at A2.
\end{grammarlens}

\subsection*{Your turn}
\begin{itemize}
\raggedright
  \item \emph{Say it:} \emph{kitaku}
  \item \emph{Say it:} \emph{kitaku}, once more
  \item \emph{Say it:} \emph{chichi no kitaku wa yūgata desu} — My father gets home in the evening.
  \item \emph{Recall:} say the Japanese for to be late, then the Japanese for going out, then say \emph{kitaku} again
\end{itemize}

\begin{tcolorbox}[breakable,colback=teal!4,colframe=teal!35!black,title={Before you move on}]
What does \ja{きたく} mean? (\textquotedblleft{}Coming home; getting home\textquotedblright{}.) Say it once more.
\end{tcolorbox}
\section[\texorpdfstring{\ja{わすれもの} (wasuremono)}{わすれもの (wasuremono)}]{\ja{わすれもの} (wasuremono) — something left behind}

\label{lesson:JA-C153-wasuremono}



\begin{quote}
Before the new one: say the Japanese for going out, then the Japanese for coming home.
\end{quote}

\subsection*{You'll want to know: \ja{わすれもの}}
\textbf{\ja{わすれもの}} — \emph{wasuremono} — \textquotedblleft{}something left behind\textquotedblright{}.

\begin{quote}
\textbf{\ja{この} \ja{かさは} \ja{ちちの} \ja{わすれものです。}} — \emph{kono kasa wa chichi no wasuremono desu} — \textquotedblleft{}This umbrella is something my father left behind.\textquotedblright{}
\end{quote}

\begin{cousinweb}
\textbf{\ja{わすれもの}} (\emph{wasuremono}) is native Japanese, built from \emph{wasure}, \textquotedblleft{}forgetting\textquotedblright{}, the stem of \emph{wasureru}, and \emph{mono}, \textquotedblleft{}thing\textquotedblright{}, both inherited rather than borrowed (JMdict, KANJIDIC2).
\end{cousinweb}

\begin{grammarlens}[title={What you've built}]
One more word for everyday life at A2.
\end{grammarlens}

\subsection*{Your turn}
\begin{itemize}
\raggedright
  \item \emph{Say it:} \emph{wasuremono}
  \item \emph{Say it:} \emph{wasuremono}, once more
  \item \emph{Say it:} \emph{kono kasa wa chichi no wasuremono desu} — This umbrella is something my father left behind.
  \item \emph{Recall:} say the Japanese for going out, then the Japanese for coming home, then say \emph{wasuremono} again
\end{itemize}

\begin{tcolorbox}[breakable,colback=teal!4,colframe=teal!35!black,title={Before you move on}]
What does \ja{わすれもの} mean? (\textquotedblleft{}Something left behind\textquotedblright{}.) Say it once more.
\end{tcolorbox}
\section[\texorpdfstring{\ja{なくす} (nakusu)}{なくす (nakusu)}]{\ja{なくす} (nakusu) — to lose; to mislay}

\label{lesson:JA-C153-nakusu}



\begin{quote}
Before the new one: say the Japanese for coming home, then the Japanese for something left behind.
\end{quote}

\subsection*{You'll want to know: \ja{なくす}}
\textbf{\ja{なくす}} — \emph{nakusu} — \textquotedblleft{}to lose; to mislay\textquotedblright{}.

\begin{quote}
\textbf{\ja{いもうとは} \ja{よく} \ja{かさを} \ja{なくします。}} — \emph{imōto wa yoku kasa o nakushimasu} — \textquotedblleft{}My younger sister often loses her umbrella.\textquotedblright{}
\end{quote}

Polite \textbf{\ja{なくします}}, \emph{nakushimasu}; a godan verb.

\begin{cousinweb}
\textbf{\ja{なくす}} (\emph{nakusu}) is a native Japanese word (wago), inherited rather than borrowed from Chinese, and its older origin is not settled. It is usually written in kana; its kanji is a later kun'yomi (JMdict, KANJIDIC2).
\end{cousinweb}

\begin{grammarlens}[title={What you've built}]
One more word for everyday life at A2.
\end{grammarlens}

\subsection*{Your turn}
\begin{itemize}
\raggedright
  \item \emph{Say it:} \emph{nakusu}
  \item \emph{Say it:} \emph{nakusu}, once more
  \item \emph{Say it:} \emph{imōto wa yoku kasa o nakushimasu} — My younger sister often loses her umbrella.
  \item \emph{Recall:} say the Japanese for coming home, then the Japanese for something left behind, then say \emph{nakusu} again
\end{itemize}

\begin{tcolorbox}[breakable,colback=teal!4,colframe=teal!35!black,title={Before you move on}]
What does \ja{なくす} mean? (\textquotedblleft{}To lose; to mislay\textquotedblright{}.) Say it once more.
\end{tcolorbox}
\section[\texorpdfstring{\ja{まにあう} (maniau)}{まにあう (maniau)}]{\ja{まにあう} (maniau) — to be in time; to make it}

\label{lesson:JA-C153-maniau}



\begin{quote}
Before the new one: say the Japanese for something left behind, then the Japanese for to lose.
\end{quote}

\subsection*{You'll want to know: \ja{まにあう}}
\textbf{\ja{まにあう}} — \emph{maniau} — \textquotedblleft{}to be in time; to make it\textquotedblright{}.

\begin{quote}
\textbf{\ja{まいあさ} \ja{でんしゃに} \ja{まにあいます。}} — \emph{maiasa densha ni maniaimasu} — \textquotedblleft{}Every morning I make it to the train in time.\textquotedblright{}
\end{quote}

Polite \textbf{\ja{まにあいます}}, \emph{maniaimasu}; a godan verb.

\begin{cousinweb}
\textbf{\ja{まにあう}} (\emph{maniau}) is native Japanese, built from \emph{ma}, \textquotedblleft{}interval, space\textquotedblright{}, the particle \emph{ni}, and \emph{au}, \textquotedblleft{}to fit\textquotedblright{}: to fit inside the time. All three parts are inherited rather than borrowed (JMdict, KANJIDIC2).
\end{cousinweb}

\begin{grammarlens}[title={What you've built}]
One more word for everyday life at A2.
\end{grammarlens}

\subsection*{Your turn}
\begin{itemize}
\raggedright
  \item \emph{Say it:} \emph{maniau}
  \item \emph{Say it:} \emph{maniau}, once more
  \item \emph{Say it:} \emph{maiasa densha ni maniaimasu} — Every morning I make it to the train in time.
  \item \emph{Recall:} say the Japanese for something left behind, then the Japanese for to lose, then say \emph{maniau} again
\end{itemize}

\begin{tcolorbox}[breakable,colback=teal!4,colframe=teal!35!black,title={Before you move on}]
What does \ja{まにあう} mean? (\textquotedblleft{}To be in time; to make it\textquotedblright{}.) Say it once more.
\end{tcolorbox}
